\documentclass[11pt]{article}
\usepackage[a4paper,margin=24mm]{geometry}
\usepackage{amsmath,amssymb,mathtools}
\usepackage[colorlinks=true]{hyperref}
\title{Two independent bubble subtractions and a triangle\\Exact three-dimensional reduction}
\author{Phi3Lab continuation}
\date{25 October 2026}
\begin{document}\maketitle

\paragraph{Original graphs.}
The cohort is G204, G253, G411 and G465 in the four-loop three-point ledger. Each graph has two disjoint one-loop two-point children and one disjoint one-loop triangle.
All eight forests are reconstructed from original edge routings.
The Gaussian quadratic matrix, positive polynomials, and Schur complement constructions of $U,F$ agree exactly.
Each local bubble retains $T^2$; the triangle retains $T^0$.
The overall $T^0$ acts on physical external momenta at fixed remaining loop momentum $q$.
The normalization is $d^6q/\pi^3$; symmetry factors and coupling are outside this integral.

\paragraph{Mass and kinetic subtraction.}
For $t\geq0$ define
\[
B(t)=\int_0^1\!\left[(1+x(1-x)t)\log(1+x(1-x)t)-x(1-x)t\right]dx .
\]
The positive sign follows from the complete local two-point Schwinger remainder.
The identities $B(0)=B'(0)=0$ track the mass and kinetic subtraction independently.
With
\[
E(t)=2\sqrt{\frac{t+4}{t}}\operatorname{arsinh}\frac{\sqrt t}{2}-2,\qquad
j_0(t)=\frac{4\operatorname{arsinh}(\sqrt t/2)}{\sqrt{t(t+4)}},
\]
the exact scalar and its derivatives are
\begin{align*}
B(t)&=\frac{t+4}{6}E(t)-\frac t9,\\
B'(t)&=\frac{E(t)}6+\frac{(t+4)(1-j_0(t))}{6t}-\frac19,\\
B''(t)&=\frac1{6t}-\frac{1-j_0(t)}{t^2}.
\end{align*}
Each bubble absorbs one \emph{retained} host propagator carrying its exact boundary momentum:
\[
J(t)=\frac{B(t)}{1+t},\quad
J'=\frac{B'}{1+t}-\frac B{(1+t)^2},\quad
J''=\frac{B''}{1+t}-\frac{2B'}{(1+t)^2}+\frac{2B}{(1+t)^3}.
\]
If two boundary momenta coincide, two different propagators are still retained.
The original graph has five host propagators before this regrouping and three afterwards.

\paragraph{Complete three-dimensional density.}
The triangle factor is $T=-L(a,b,1)$, including simplex volume $1/2$.
For the physical scaling parameter $\lambda$, let
\[
f(\lambda)=H(\lambda)J(t_1(\lambda))J(t_2(\lambda))T(\lambda),
\qquad
H(\lambda)=\prod_{e\in\text{three remaining hosts}}
\frac{1+t}{1+|\epsilon_e q+\lambda b_e|^2}.
\]
At zero physical momentum, $f(0)=J(t)^2[-E(t)/2]$, with $t=q^2$ retained.
The full forest sum is represented by
\[
\frac{f(1)+f(-1)}2-f(0).
\]
No constituent forest is discarded.
Set $t=(s/(1-s))^2$, with $s,z,\theta\in(0,1)$.
The full density is the antipodal remainder times
\[
\frac{2}{s(1-s)}\left(\frac{t}{1+t}\right)^3(1-z).
\]
The radial measure after angular normalization is
$t^2\,dt/[2(1+t)^3]$.

\paragraph{Quadratic control and factorial.}
For a routed bubble invariant $t_b(\lambda)=t+\lambda\ell+\lambda^2c$,
its Taylor coefficients through degree two are
\[
J(t),\quad J'(t)\ell,\quad J'(t)c+\tfrac12J''(t)\ell^2.
\]
The coefficient $\ell=2\epsilon q\cdot b$ uses the complete original boundary momentum.
All three host ratios, both bubble factors and the triangle are convolved before taking the quadratic coefficient $C=f''(0)/2$.
Subtract $C(q)$ and return its exact six-dimensional angular mean:
\[
\langle(q\cdot p_a)(q\cdot p_b)\rangle=\frac{t}{6}(p_a\cdot p_b).
\]
In cube coordinates the added density is
\[
\frac{2}{s(1-s)}\left(\frac{t}{1+t}\right)^3
\left[\frac{\langle C\rangle}{2}-(1-z)C(q)\right].
\]
Its integral is exactly zero. It is an integration control, not an additional physical counterterm.

\paragraph{Independent representation.}
The independent reference keeps \emph{all five original host denominators}, uses $B$ instead of $J$, and forms the complete product before the overall subtraction.
Its common radial factor has an additional $(1+t)^{-2}$.
That reference computes Taylor coefficients by formal convolution and the angular mean by twelve axial directions; the native kernel uses an explicit covariance matrix.
These are distinct implementations of the regrouping and angular control.

\paragraph{Absolute integrability.}
For $t\geq16$, write $r=\sqrt t$, $L=\log(1+t)$.
All shifts satisfy $|b|\leq1$, so
$t/2\leq|q+\lambda b|^2\leq2t$ for $|\lambda|\leq1$.
From
\[
0\leq B(s)\leq sB'(s),\qquad
B'(s)\leq\tfrac16\log(1+s),\qquad B''(s)\leq(6s)^{-1}
\]
one obtains
\[
|J|\leq L,\quad|\partial_\lambda J|\leq2L/r,\quad
|\partial_\lambda^2J|\leq40L/t .
\]
The already proved triangle bounds are $|T|\leq L$, $|T'|\leq4/r$, $|T''|\leq16L/t$.
For $P=J_1J_2T$ this gives $|P'|\leq6L^3/r$ and $|P''|\leq120L^3/t$.
The three-host ratio satisfies
$|H|\leq8$, $|H'|\leq120/r$, $|H''|\leq2496/t$.
Consequently $|f''|\leq4896L^3/t$.
The antipodal remainder and the mean-return control yield the absolute radial bound
\[
\frac{3672\log^3(1+t)}{t^2}.
\]
For $T\geq16$ its integrated tail is bounded by
\[
\frac{3672}{T}\left[\log^3(2T)+3\log^2(2T)+6\log(2T)+6\right].
\]
On $0\leq t\leq16$, an explicit bound $5570t^2$ covers the origin and compact region.
These bounds establish absolute integrability over the whole axis; they do not certify a numerical precision target.

\paragraph{Evidence.}
The original factor proof is
\texttt{verify\_disjoint\_mixed\_three\_children\_v1.py}.
The control and whole-axis proof is
\texttt{verify\_disjoint\_mixed\_three\_control\_v1.py}.
Read actual supervisor and worker exits and their frozen source hashes.
Kernel interval checks, independent direct child integration, fixed diagnostics, independent production and full stored-request replay are required before accounting.
\end{document}
